\begin{textsample}{NEG Dim 4 – human – Score -2.00 – mg/mg\_ef\_ch\_geo\_3\_p2.txt}  \label{ex:f4_neg_007}
Nesse sentido, o componente \textbf{curricular} de Geografia, por meio do desenvolvimento de competências e habilidades, favorece a formação integral do estudante proporcionando -lhe a compreensão do espaço local e suas inter-relações com o global ; permitindo -lhe o reconhecimento e a valorização dos diferentes saberes como também favorecendo a construção de uma atitude responsável diante da natureza e consciente de suas responsabilidades e direitos.
O Currículo Referência de Minas Gerais em consonância com a Base \textbf{Nacional} Comum \textbf{Curricular} tem também como princípio a inclusão da diversidade mais ampla de todos os estudantes.

% matched lemmas: curricular, nacional
\end{textsample}
